\section{Electronic Noise}
\subsection{Presentation}

In the acoustical domain, noise is usually associated with undesired disturbance. In the electronical domain, it may be desired.  

The standard notation for noise is:
\[
v(t) = V_{DC} + v_s(t) + v_n(t) + v_d(t)
\]
Where:
\begin{itemize}
    \item $v_n(t)$ : the noise voltage
    \item $v_d(t)$ : the disturbance voltage (for example the 50 Hz AC supply)
\end{itemize}

\textbf{NB:} Random signals are often associated to zero mean value, i.e. $ \overline{v_n} = 0 $.

The average power is usually computed taking the RMS value of the signals:
\[
v_{n,RMS} = \sqrt{\lim_{T \to +\infty} \frac{1}{T} \int_0^T v_n^2(t)\, dt}
\]

We introduce:
\begin{itemize}
    \item the \textbf{Signal to Noise Ratio (SNR)}: 
    \[
    SNR = \frac{v_{s,RMS}}{v_{n,RMS}} \in [0,+\infty[
    \]
    \item the \textbf{Signal to Disturbance Ratio (SDR)}:
    \[
    SDR = \frac{v_{s,RMS}}{v_{d,RMS}} \in [0,+\infty[
    \]
    \item Both quantities are usually expressed in dB:
\[
SNR_{dB} = 20 \cdot \log_{10}(SNR)
\]
\end{itemize}

$v_{n,RMS}(f_0, \Delta f)$ is characterized by the central frequency $f_0$ and its bandwidth $\Delta f$.

\textbf{NB:} SNR at low frequencies is lower than SNR at high frequencies.  

\begin{figure}[H]
    \centering
    \includegraphics[width=0.7\linewidth]{Narrow_Band_Filtering.png}
    \caption{Processing pipeline for spectral distribution analysis.}
    \label{fig:placeholder}
\end{figure}

The \textbf{Noise Power Spectral Density (NPSD)} plot describes how the squared RMS voltage behaves on the frequency spectrum.

We identify different components in the Noise Power Spectral Density (NPSD) curve:
\begin{itemize}
    \item \textbf{White noise component} $S_{n,w}(v^2)$: constant across the spectrum.
    \item \textbf{Flicker / Pink noise component} $S_{n,f}(v^2)$: decreases with frequency in a $\propto \tfrac{c}{f}$ manner.
    \item \textbf{Brown noise component} $S_{n,b}(v^2)$: decreases faster, in a $\propto \tfrac{c}{f^2}$ manner.
\end{itemize}

\textbf{Remark:} Pink noise is particularly relevant in audio systems, as it impacts the low-frequency part of the spectrum.  
The typical audio band (20--20 kHz) is often considered as a \textit{low-frequency range} in electronic circuits.  

When analyzing noise, we can choose voltage or current as our reference for Power Spectral Density (PSD):
\[
S_{v,f} = \frac{dv_n^2}{df} \qquad [V^2/Hz]
, \qquad
S_{i,f} = \frac{di_n^2}{df} \qquad [A^2/Hz]
\]
Thus, we usually refer to \textbf{Voltage Noise Spectral Density} $\left( S_{v,f}\right)$ and \textbf{Current Noise Spectral Density} $\left( S_{i,f}\right)$.  

By integrating the PSD over frequency, we obtain the RMS values of these noise components:
\[
\int_{BW} S_n(f)\, df = \overline{v_n^2} = v_{RMS,n}^2
\quad \text{or} \quad
\int_{BW} S_n(f)\, df = \overline{i_n^2} = i_{RMS,n}^2
\]
depending on whether the target is voltage or current.

\subsection{Measuring Noise}

\begin{figure}[H]
    \centering
    % Left side: explanation + equations
    \begin{minipage}[c]{0.6 \linewidth}
    To evaluate the overall contribution of each component (voltage and current) to the output, we resort to an \textbf{equivalent circuit model}.  
Random noise components may add together or cancel out; therefore, their relationship is often analyzed in terms of correlation.  

We typically use \textbf{two generators} at the input, instead of one, for the following reasons:
    \end{minipage}\hfill
    \begin{minipage}[c]{0.3\linewidth}
        \centering
        \includegraphics[width=\linewidth]{Meausing_Noise_1.png}
    \caption{Abstract representation of the setup.}
    \end{minipage}
\end{figure}

\begin{itemize}
    \item If we measure the output voltage with \textbf{no input voltage}, the noise voltage source contributes nothing (since it is floating).  
    The measured noise is then solely due to the \textbf{current generator}.
    
    \item If we \textbf{short-circuit the input}, the current generator no longer contributes.  
    The measured output noise is then entirely due to the \textbf{voltage generator}.
\end{itemize}

\textbf{Conclusion:}  
By combining these two measurement cases, we can separately identify and quantify the contributions of voltage noise and current noise to the total system noise.

We know that for any circuit, the output is given by :
\[
{{v}_{out,sig}}={{v}_{in,sig}}\cdot TF\left( j\omega  \right)
\]
When working with the NPSD $S_{out,n} \ [V^2/Hz]$, such that $S_{out,n} = | TF(j \omega) |^2 \cdot S_{in,n}$.


\subsection{Resistors' Noise}

We can think of resistors as blocks of conductors to which we apply a certain voltage and current.  
Inside the conductor, electric charges move randomly due to thermal agitation.  
This random motion generates fluctuations, which create heat in the conductor.  
This phenomenon is called \textbf{thermal noise} or \textbf{Johnson noise}.

The Voltage Noise PSD of a Resistor is given by: 
\[
S^R_{n,V} = 4k_B T R
\]
which is \textbf{frequency independent}.  

\begin{itemize}
    \item $k_B = 1.38 \cdot 10^{-23}\ [J/K]$ is the Boltzmann constant.
    \item $T$ is the absolute temperature in Kelvin (commonly $T \approx 300\,K$).
    \item $R$ is the resistance in $\Omega$.
\end{itemize}

Since $S^R_{n,v}$ does not depend on frequency, this is a \textbf{white noise process}.

The RMS fluctuation is computed as:
\[
v_{n,RMS} = \sqrt{\int_{BW} S_{n,V} \, df}
= \sqrt{\int_{BW} 4 k_B T R \, df}
= \sqrt{BW \cdot 4k_B T R}
\]
For an audio bandwidth of $20\,\text{Hz} - 20\,\text{kHz}$:
\[
v_{n,RMS} = \sqrt{(20 \cdot 10^3 - 20)\, 4k_B T R}
\approx \sqrt{20 \cdot 10^3 \cdot 4k_B T R}
\]

Numerical Example: taking $R = 1\,k\Omega$ and $T = 300\,K$:
\[
S_{n,v} \approx 1.6 \cdot 10^{-17}\ \left[\frac{V^2}{Hz}\right], 
\qquad v_{n,RMS} \approx 80 \,\mu V_{\text{RMS}}
\]
From Ohm’s law $V = RI$, the current noise PSD is obtained as:
\[
S_{n,i} = \frac{S_{n,v}}{|R|^2} 
= \frac{4k_B T R}{R^2}
\quad \Rightarrow \quad
S_{n,i} = \frac{4k_B T}{R}
\]

\subsection{PN Junctions Noise}

The noise described in the figure is called \textbf{shot noise}.  
It is associated with the current flow, more precisely related to the passage of electrons from one side to the other of the PN junction.  
\begin{figure}[H]
    \centering
    \includegraphics[width=0.8\linewidth]{shot_noise.png}
    \caption{Illustration of the shot noise.}
\end{figure}
\[
S_{n,I} = \frac{d\overline{i_n}^2}{df} = 2 \cdot q \cdot I
\]
where $q = 1.6 \cdot 10^{-19}$ C is the electron charge, and $I$ is the current through the junction. We observe that this is a \textbf{white noise} process, i.e. constant in frequency.  

\paragraph{Comparison: Thermal vs Shot Noise}
\begin{center}
\begin{tabular}{|c|c|}
\hline
\textbf{Thermal noise} & \textbf{Shot noise} \\
\hline
Depends on the resistor (material, device) & Depends on the current (not the material, but the bias) \\
\hline
Depends on temperature (can be reduced) & Cannot be reduced by temperature \\
\hline
\end{tabular}
\end{center}

\clearpage

\subsection{Noise in Transistors}

We have different transistor types: BJT, JFET, and MOSFET.

\begin{figure}[H]
    \centering
    % Left side: explanation + equations
    \begin{minipage}[c]{0.75 \linewidth}
    \paragraph{BJT:}  
The main noise component is a current noise seen across its emitter and collector, $\overline{i_{n,C}^2}$, associated to the CNPSD:
\[
S_{n,I} = 2q \cdot I_C \text{ which is purely shot noise.}
\]
The Voltage Noise PSD of a BJT is:
\[
S_{n,V}^{BJT} = 2q \cdot \frac{I_C}{g_m^2} \text{ where $g_m$ is the transconductance.}
\]
    \end{minipage}\hfill
    \begin{minipage}[c]{0.2\linewidth}
        \centering
        \includegraphics[width=\linewidth]{BJT_Noise.png}
    \caption{Model of the noise within a BJT.}
    \end{minipage}\hfill
\end{figure}

\begin{figure}[H]
    \centering
    % Left side: explanation + equations
    \begin{minipage}[c]{0.75 \linewidth}
    \paragraph{JFETs and MOSFETs:}  
The current noise is associated with the current flow through the gate $\overline{i_{n,G}^2}$, with CNPSD:
\[
S_{n,G}^{MOSFET/JFET} = 2q I_G \ \ \text{ where $I_G$ is very small (pA for JFETs, fA for MOSFETs).}
\]
Another important source of noise is the \textbf{drain current noise}, $\overline{i_{n,D}^2}$, modeled as thermal noise:
\[
S_{n,I}^{MOSFET/JFET} = \frac{4 k_B T}{R_{eq}} = 4k_B T \cdot g_m
\quad \text{where } g_m = \frac{1}{R_{eq}}
\]
    \end{minipage}\hfill
    \begin{minipage}[c]{0.2\linewidth}
        \centering
        \includegraphics[width=\linewidth]{JFET_MOSFET_Noise.png}
    \caption{Model of the noise within a JFET/MOSFET.}
    \end{minipage}\hfill
\end{figure}

Thus, the Voltage Noise PSD of JFETs and MOSFETs is:
\[
S_{n,V}^{JFET/MOSFET} = \frac{4k_BT g_m}{g_m^2} = \frac{4k_BT}{g_m}
\]

\paragraph{Corner Frequency:}
It is the frequency where the two asymptotic noise curves merge.

\begin{itemize}
    \item For JFETs and BJTs, the corner frequency is typically around $10\,Hz$.
    \item For MOSFETs, the corner frequency is much higher, often above the audio range.
\end{itemize}

\textbf{Remark:}  
Since MOSFET corner frequencies are so high, they are generally not used in audio systems.  
However, MOSFETs remain dominant in integrated circuits because of their high integration capability (smaller size than other transistor types).

\begin{figure}[H]
    \centering
    \includegraphics[width=0.8\linewidth]{Corner_Frequency.png}
    \caption{NPSD of transistors with the corner frequency.}
    \label{fig:placeholder}
\end{figure}

\clearpage


\subsection{Noise in Multi-Stage Systems}

Let’s consider a 3-stage circuit:
\begin{figure}[H]
    \centering
    \includegraphics[width=0.8\linewidth]{Multi_Stage_Noise_Pipeline.png}
    \caption{Pipeline of 3 OPAMPs with noise.}
\end{figure}

\textbf{NB:} The red and blue quantities indicate noise sources.  

The output signal is the product of the input signal and the gain stages.  
The output \textit{noise}, however, is different: to compute it, we must account for the \textbf{impedances in the circuit} as illustrated below.
\begin{figure}[H]
    \centering
    \includegraphics[width=0.25\linewidth]{Source_Impedane_Noise.png}
    \caption{Accounting for the impedances at stage 1.}
\end{figure}

\begin{figure}[H]
    \centering
    % Left side: explanation + equations
    \begin{minipage}[c]{0.55 \linewidth}
    \textbf{Current noise component in the 1\textsuperscript{st} stage:}
We remove (short-circuit) the voltage source and the voltage noise source.  
The voltage associated to the input current is:
\[
\overline{v_{n,1_{IN}}}^2 = \overline{i_{n,1}}^2 \cdot \left|Z_{in,1}//Z_S\right|^2
\quad \Rightarrow \quad
\overline{v_{n,1_0}}^2 = \overline{i_{n,1}}^2 \cdot \left|Z_{in,1}//Z_S\right|^2 \cdot \left|G_1\right|^2
\]
    \end{minipage}\hfill
    \begin{minipage}[c]{0.35\linewidth}
        \centering
        \includegraphics[width=\linewidth]{Current_Component.png}
    \caption{Circuit with the voltage sources as short circuits.}
    \end{minipage}\hfill
\end{figure}


\begin{figure}[H]
    \centering
    % Left side: explanation + equations
    \begin{minipage}[c]{0.55 \linewidth}
    \textbf{Voltage noise component in the 1\textsuperscript{st} stage:}
We remove (open-circuit) the current source and the current noise source.  
The current associated to the input voltage is:
\[
\overline{v_{n,1_{IN}}}^2 = \overline{v_{n,1}}^2 \cdot \left|\frac{Z_{in,1}}{Z_{in,1} + Z_s}\right|^2
\quad \Rightarrow \quad
\overline{v_{n,1_0}}^2 = \overline{v_{n,1}} ^2 \cdot \left|\frac{Z_{in,1}}{Z_{in,1} + Z_s}\right|^2 \cdot |G_1|^2
\]
    \end{minipage}\hfill
    \begin{minipage}[c]{0.35\linewidth}
        \centering
        \includegraphics[width=\linewidth]{Voltage_Component.png}
    \caption{Circuit with the current sources as floating points.}
    \end{minipage}\hfill
\end{figure}

\textbf{Complete Output Noise:}
The total output noise contribution of stage 1 to the output of stage 1 is:
\[
\overline{v_{n,1_0}}^2 = \overline{i_{n,1}}^2 \left|Z_{in,1}//Z_S\right|^2 \cdot |G_1|^2 
+ \overline{v_{n,1}}^2 \cdot \left|\frac{Z_{in,1}}{Z_{in,1} + Z_s}\right|^2 \cdot |G_1|^2
\]
\textbf{Discussion:}
\begin{itemize}
    \item The \textbf{current noise} becomes more relevant as the source impedance increases.  
    This is one of the reasons why we typically design systems with \textbf{low source impedances}.
    \item When $Z_s \to 0$, the \textbf{voltage noise} dominates.  
    \item When $Z_s \to \infty$, the \textbf{current noise} dominates.
\end{itemize}


\begin{figure}[H]
    \centering
    \includegraphics[width=0.8\linewidth]{Multi_Stage_Noise_params.png}
    \caption{Pipeline of the multi-stage system with $\alpha, \beta, \gamma$} parameters.
\end{figure}
\textbf{NB:} Parameters $\alpha, \beta, \gamma$ account for the impedances in parallel at each stage of the pipeline.


Each input is given by a voltage partition:
\[
\alpha =\frac{{{Z}_{IN,1}}}{{{Z}_{IN,1}}+{{Z}_{S}}},\ \beta =\frac{{{Z}_{IN,2}}}{{{Z}_{IN,2}}+{{Z}_{0,1}}},\ \gamma =\frac{{{Z}_{IN,3}}}{{{Z}_{IN,3}}+{{Z}_{0,2}}}
\]
and we have:
\[
v_{out,s} = \alpha G_1 \cdot \beta G_2 \cdot \gamma G_3 \cdot v_s
\]

\textbf{Noise Voltage Contribution of Stage 1: }
The noise voltage of stage 1 is filtered by both $TF_2$ and $TF_3$.  
Thus, the contribution of stage 1 noise ($\overline{v}_{out,n,1,0}^2$) on the system output is:
\[
\overline{v}_{out,N,1}^2 = \overline{v}_{n,1,0}^2 \cdot |TF_2 \cdot TF_3|^2 
= \overline{v}_{n,1,0}^2 \cdot |\beta G_2 \cdot \gamma G_3|^2
\]
The noise PSD at input ($S_{N,1}$) and output ($S_{N,1_0}$)of stage 1 are:
\[
S_{N,1} = \overline{v_{n,1}}^2 + \overline{v_{n,1_{IN}}}^2  
\quad \Rightarrow \quad
S_{N,1_0} = S_{N,1}\,|TF_1|^2 = S_{N,1} \cdot |\alpha G_1|^2
\]
The noise PSD contribution of stage 1 in the entire system is:
\[
S_{OUT,N1} = S_{N,1}\,|TF_1  \cdot TF_2 \cdot TF_3|^2
\]
\textbf{Stage 2 and Stage 3 Contributions: }
Similarly to stage 1:
\[
S_{OUT,N,2} = S_{N,2}\,|TF_2 \cdot TF_3|^2, 
\qquad
S_{OUT,N,3} = S_{N,3} \left| TF_3 \right|^2
\]
The total output noise PSD is the sum of the three contributions:
\[
S_{OUT,N} = S_{OUT,N,1} + S_{OUT,N,2} + S_{OUT,N,3}
= S_{N,1}\,|TF_1 \cdot TF_2 \cdot TF_3|^2 + S_{N,2}\,|TF_2 \cdot TF_3|^2 + S_{N,3}|TF_3|^2
\]
The noisiness of the system is expressed by the signal-to-noise ratio:
\[
SNR = \frac{v_{OUT,S,RMS}}{v_{OUT,N,RMS}}
= \frac{v_{S,RMS}\,|TF_1 \cdot TF_2 \cdot TF_3|}
{\sqrt{\int_{BW} S_{OUT,N}\,df}}
= \frac{v_{S,RMS}}{\sqrt{\int_{BW} \left(S_{N,1} + \frac{S_{N,2}}{|TF_1|^2} + \frac{S_{N,3}}{|TF_1 \cdot TF_2|^2}\right)df}}
\]
\textbf{Discussion}
\begin{itemize}
    \item The noise generated at the \textbf{first stage} contributes directly to the output SNR.  
    \item Noise from later stages is attenuated when referred back to the input, because it is divided by the gains of previous stages.  
    \item Therefore, the \textbf{first stage} plays the most critical role in determining the system’s noise performance.
\end{itemize}

\textbf{NB:} For audio applications, the bandwidth is typically:
\[
BW = [20, \, 20 \cdot 10^3] \ \text{Hz}
\]
\textbf{Key Design Principles for Low Noise Operation}
\begin{itemize}
    \item The first stage should have as little noise as possible.  
    \item The first stage should provide as much gain as possible.  
\end{itemize}

This is why the \textbf{preamplifier} is such a crucial component in audio systems.


\subsection{Noise in OPAMPs}

An OPAMP disconnected from any external signal has two main internal noise sources: a voltage noise source and two bias current noise sources.
\begin{figure}[H]
    \centering
    \includegraphics[width=0.25\linewidth]{Noise_OPAMP.png}
    \caption{Noise model in the OPAMP.}
\end{figure}



\subsubsection{Exercise: Noise in a Non-Inverting Configuration}


\begin{figure}[H]
    \centering
    % Left side: explanation + equations
    \begin{minipage}[c]{0.6 \linewidth}
    For a non-inverting amplifier:
\[
v_{out} = \left(1 + \frac{R_2}{R_1}\right) v_s, 
\qquad G = 1 + \frac{R_2}{R_1}
\]
The noise components arise from the resistors:
\[
\quad \overline{v_{n,Rs}}^2, \quad \overline{v_{n,R1}}^2, \quad \overline{v_{n,R2}}^2
\]
The NPSD of a generic resistor is:
\[
S_{n,R} = 4k_B T R
\]
    \end{minipage}\hfill
    \begin{minipage}[c]{0.4\linewidth}
        \centering
        \includegraphics[width=\linewidth]{Noise_OPAMP_Ex.png}
    \caption{Non-inverting OPAMP of interest.}
    \end{minipage}\hfill
\end{figure}
Thusthe output NPSD is:
\[
S_{out,n,Rx} = |TF_{Rx}|^2 \cdot S_{n,Rx}
\]
\textbf{Contribution of $R_s$:}
\[
S_{out,n,Rs} = |G|^2 \cdot S_{n,Rs} = |G|^2 4 \cdot k_b \cdot T \cdot R_s
\]
\textbf{Contribution of $R_1$:}
\[
S_{out,n,R1} = |G-1|^2 \cdot 4 \cdot k_B \cdot T \cdot R_1
\]
\textbf{Contribution of $R_2$:}
We must reason carefully:  
If the voltage sources are short-circuited, there is no voltage across $R_2$.  
Since the negative input of the OPAMP is at ground potential (virtual ground), we also have 0V across $R_2$.  
Thus the noise contribution of $R_2$ is:
\[
S_{out,n,R2} = 4\cdot k_B \cdot T \cdot R_2
\]

\begin{figure}[H]
    \centering
    \includegraphics[width=0.6\linewidth]{noise_volt_curr_noise.png}
    \caption{Scheme of the OPAMP with voltage noise $\overline{e}_n^2$ and input current noises $\overline{i}_{n_+}^2, \ \overline{i}_{n_-}^2$.}
\end{figure}
\textbf{Other Contributions:}
\begin{itemize}
    \item Input voltage noise $\overline{e_n}^2$: 
    \[
    S_{out,n,e} = |G|^2 \cdot S_{n,e}, \quad \text{(given in datasheets)}
    \]
    \item Input current noise $\overline{i}_{n_+}^2$ (through $R_s$):
    \[
    S_{out,n,i_+} = |R_s \cdot G|^2 \cdot S_{n,i_+}
    \]
    \item Input current noise $\overline{i}_{n_-}^2$ (through $R_2$):  
    Reasoning as before, the current noise must flow through $R_2$ (with negative input output voltage at $R_2$).  
    Thus:
    \[
    S_{out,n,i_-} = |-R_2|^2 \cdot S_{n,i_-}
    \]
\end{itemize}

The \textbf{total Output Noise Power Spectral Density is:}
\[
S_{out,n,tot} = S_{out,n,Rs} + S_{out,n,R1} + S_{out,n,R2} + S_{out,n,e} + S_{out,n,i_{-}} + S_{out,n,i_{+}}
\]
The \textbf{RMS Output Noise is} : 
\[
v_{out,n,RMS} = \sqrt{\int_{BW} S_{out,n,tot}\, df}
\]
If $S_{out,n,tot}$ is white (i.e. constant across frequency), then:
\[
v_{out,n,RMS} = \sqrt{(f_h - f_l) \cdot S_{out,n,tot}}
\]
For the audio band with $ BW = [20,\, 20 \cdot 10^3] \ \text{Hz} $
\[
\Rightarrow \quad v_{out,n,RMS} \approx \sqrt{20 \cdot 10^3} \cdot \sqrt{S_{out,n,tot}}
\]
The \textbf{Signal-to-Noise Ratio (SNR)} is: 
\[
SNR = \frac{v_{out,s,RMS}}{v_{out,n,RMS}} = \frac{G \cdot v_s}{v_{out,n,RMS}}
\]


\subsection{Notes from the slide}
In audio and electroacoustic systems, the \textbf{noise performance} of electronic components—especially operational amplifiers (op-amps)—is crucial in determining the achievable signal-to-noise ratio (SNR) and dynamic range.  
Datasheets provide detailed noise specifications to help engineers predict and minimize noise in preamplifiers, filters, and other low-level audio circuits.

This section summarizes how noise is characterized, how to read it from datasheets, and how it affects the performance of practical audio systems.

\paragraph{Types of Electronic Noise}
The main types of intrinsic noise encountered in electronic components are:
\begin{itemize}
  \item \textbf{Thermal noise (Johnson–Nyquist noise):} generated by random motion of charge carriers in resistors and semiconductors.  
        It is white noise with spectral density:
        \[
        e_n = \sqrt{4 k T R}\,[\mathrm{V}/\sqrt{\mathrm{Hz}}],
        \]
        where \(k\) is Boltzmann’s constant, \(T\) the absolute temperature, and \(R\) the resistance.
  \item \textbf{Shot noise:} produced by the discrete nature of electric charge in active devices (transistors, diodes).
  \item \textbf{Flicker noise (1/f noise):} dominant at low frequencies (below a few hundred Hz), often visible as a rising slope in log–log spectral plots.
\end{itemize}

\paragraph{Noise Parameters in Datasheets}
Operational amplifier datasheets specify several key noise parameters:
\begin{itemize}
    \item \textbf{Input Voltage Noise Density (\(e_n\)):} The equivalent input voltage noise per unit bandwidth, typically expressed in $\mathrm{nV}/\sqrt{\mathrm{Hz}}$.  
For example, a low-noise audio op-amp might specify \(e_n = 3\,\mathrm{nV}/\sqrt{\mathrm{Hz}}\) at 1\,kHz.
\item \textbf{Input Current Noise Density (\(i_n\)):}
The equivalent input current noise per unit bandwidth, expressed in $\mathrm{pA}/\sqrt{\mathrm{Hz}}$.  
It becomes significant when the source resistance is high, since voltage noise due to current noise is \(V = i_n R_S\).
\item \textbf{Noise Spectral Density Graph:}
Datasheets usually plot $e_n(f)$ or total input noise vs. frequency.  
At low frequencies, flicker noise increases, while above a few hundred Hz, the noise tends to flatten (white region).
\item \textbf{4. Total Noise Voltage:}
The total RMS noise over a given bandwidth \(B\) is obtained by integrating the noise spectral density:
\[
V_{n,\text{RMS}} = \sqrt{ \int_{0}^{B} e_n^2(f) \, df }.
\]
For white noise, this simplifies to \(V_{n,\text{RMS}} = e_n \sqrt{B}\).
\item \textbf{5. Equivalent Noise Resistance:}
Some datasheets specify the \textbf{equivalent noise resistance} \(R_n\), defined such that its thermal noise equals the amplifier’s input noise:
\[
e_n = \sqrt{4 k T R_n}.
\]
This provides a convenient measure of noise performance independent of circuit details.

\end{itemize}
\begin{figure}[H]
    \centering
    \includegraphics[width=0.5\linewidth]{Input_Voltage_Noise_Density.png}
    \caption{Typical input voltage noise density of an audio op-amp.  
The low-frequency region shows 1/f noise, while above 1\,kHz the curve becomes flat (white noise).}
\end{figure}


Noise datasheets usually include the following graphs:
\begin{itemize}
  \item \textbf{Voltage noise spectral density:} $e_n$ vs. frequency (log–log scale).
  \item \textbf{Current noise spectral density:} $i_n$ vs. frequency.
  \item \textbf{Noise vs. source resistance:} showing the optimum source impedance that minimizes total noise.
\end{itemize}

The total input noise voltage $V_{n,\text{tot}}$ for a given source resistance \(R_S\) combines both effects:
\[
V_{n,\text{tot}} = \sqrt{ (e_n)^2 + (i_n R_S)^2 + 4 k T R_S }.
\]
Minimizing \(V_{n,\text{tot}}\) involves matching the amplifier’s input noise characteristics with the expected source impedance.

\textbf{Typical Values for Audio Op-Amps: }
\begin{center}
\begin{tabular}{|l|c|c|c|}
\hline
\textbf{Op-Amp Model} & \textbf{$e_n$ (nV/$\sqrt{\text{Hz}}$)} & \textbf{$i_n$ (pA/$\sqrt{\text{Hz}}$)} & \textbf{Comments} \\
\hline
NE5532 & 5 & 0.7 & Industry standard for audio preamps \\
OPA2134 & 8 & 1.0 & Low noise, FET input \\
LM4562 & 2.7 & 0.4 & Ultra-low-noise audio op-amp \\
TL072 & 18 & 2.0 & Higher voltage noise, low current noise (FET) \\
\hline
\end{tabular}
\end{center}

\paragraph{Noise in Multi-Stage Systems}
In multi-stage amplifiers (e.g., microphone preamplifier $\rightarrow$ tone control $\rightarrow$ power amplifier), the first stage typically dominates overall noise performance because its noise is amplified by all subsequent stages.

The \textbf{Friis formula} expresses total input-referred noise in cascaded stages:
\[
F_{\text{total}} = F_1 + \frac{F_2 - 1}{G_1} + \frac{F_3 - 1}{G_1 G_2} + \cdots
\]
where \(F_i\) is the noise factor and \(G_i\) the power gain of stage \(i\).

Thus, optimizing the first amplifier stage (preamplifier) is the most effective way to improve total system SNR.

\subsection{Datasheet Example Interpretation and Practical Comparison}

To understand how datasheet parameters translate into real audio performance, let us analyze two common low-noise operational amplifiers widely used in sound engineering:
\begin{itemize}
  \item \textbf{NE5532:} a bipolar-input dual op-amp, optimized for low voltage noise and high current drive capability.
  \item \textbf{OPA2134:} a JFET-input dual op-amp, designed for very low distortion and high input impedance.
\end{itemize}

\subsubsection{Theoretical Comparison Based on Datasheet Parameters}
\begin{center}
\begin{tabular}{|l|c|c|c|}
\hline
\textbf{Parameter} & \textbf{Symbol} & \textbf{NE5532} & \textbf{OPA2134} \\
\hline
Input voltage noise density & $e_n$ & $5\,\mathrm{nV}/\sqrt{\mathrm{Hz}}$ & $8\,\mathrm{nV}/\sqrt{\mathrm{Hz}}$ \\
Input current noise density & $i_n$ & $0.7\,\mathrm{pA}/\sqrt{\mathrm{Hz}}$ & $1.0\,\mathrm{fA}/\sqrt{\mathrm{Hz}}$ \\
Input type &  & Bipolar (BJT) & JFET \\
Typical input resistance &  & $100\,\text{k}\Omega$ & $10\,\text{M}\Omega$ \\
\hline
\end{tabular}
\end{center}

Both devices exhibit excellent audio performance, but their optimal applications differ depending on the source impedance:
\begin{itemize}
  \item The NE5532, with its very low voltage noise, excels with \emph{low-impedance sources} (e.g., microphone preamps, line-level circuits).
  \item The OPA2134, with extremely low current noise, is ideal for \emph{high-impedance sources} (e.g., piezo pickups, sensor inputs).
\end{itemize}

\subsubsection{Total Input-Referred Noise Calculation}

The total equivalent input noise is given by:
\[
V_{n,\text{tot}} = \sqrt{(e_n)^2 + (i_n R_S)^2 + 4kTR_S}.
\]

Assume a bandwidth of $B = 20\,\text{kHz}$ and ambient temperature $T = 300\,\text{K}$.

\paragraph{Example A — NE5532 with $R_S = 1\,\text{k}\Omega$:}
\[
\begin{aligned}
e_n &= 5\,\mathrm{nV}/\sqrt{\mathrm{Hz}}, \\
i_n R_S &= 0.7\,\mathrm{pA}/\sqrt{\mathrm{Hz}} \times 1\,\text{k}\Omega = 0.7\,\mathrm{nV}/\sqrt{\mathrm{Hz}}, \\
\sqrt{4kTR_S} &= \sqrt{4 \times 1.38\times10^{-23} \times 300 \times 10^3} = 4.1\,\mathrm{nV}/\sqrt{\mathrm{Hz}}.
\end{aligned}
\]

Thus:
\[
V_{n,\text{tot}} = \sqrt{5^2 + 0.7^2 + 4.1^2} = 6.5\,\mathrm{nV}/\sqrt{\mathrm{Hz}}.
\]
Over the 20\,kHz audio band:
\[
V_{n,\text{RMS}} = 6.5 \times \sqrt{20{,}000} = 0.92\,\mu\text{V}_{\mathrm{RMS}}.
\]

\paragraph{Example B — OPA2134 with $R_S = 100\,\text{k}\Omega$:}
\[
\begin{aligned}
e_n &= 8\,\mathrm{nV}/\sqrt{\mathrm{Hz}}, \\
i_n R_S &= 1\,\mathrm{fA}/\sqrt{\mathrm{Hz}} \times 100\,\text{k}\Omega = 0.1\,\mathrm{nV}/\sqrt{\mathrm{Hz}}, \\
\sqrt{4kTR_S} &= \sqrt{4 \times 1.38\times10^{-23} \times 300 \times 10^5} = 12.9\,\mathrm{nV}/\sqrt{\mathrm{Hz}}.
\end{aligned}
\]

Thus:
\[
V_{n,\text{tot}} = \sqrt{8^2 + 0.1^2 + 12.9^2} = 15.2\,\mathrm{nV}/\sqrt{\mathrm{Hz}},
\]
and over 20\,kHz:
\[
V_{n,\text{RMS}} = 15.2 \times \sqrt{20{,}000} = 2.15\,\mu\text{V}_{\mathrm{RMS}}.
\]

\subsubsection{Discussion of Results}

\begin{itemize}
  \item For low source impedances ($R_S < 2\,\text{k}\Omega$), the NE5532 achieves sub-microvolt noise levels, making it ideal for microphone preamplifiers, active filters, and summing stages.
  \item For high source impedances ($R_S > 50\,\text{k}\Omega$), the OPA2134 avoids the high current noise typical of bipolar op-amps and offers quieter operation despite higher voltage noise.
  \item The thermal noise of the source resistor often dominates once $R_S$ exceeds a few tens of kilohms.
\end{itemize}

\textbf{Summary of the Comparison:}
\begin{itemize}
  \item \textbf{NE5532:} lower noise for low-$R_S$ sources, excellent for professional balanced microphone inputs.
  \item \textbf{OPA2134:} superior for high-$R_S$ sources, high input impedance, and extremely low distortion.
  \item Both devices, properly applied, provide noise levels well below the audibility threshold in typical studio or live sound applications.
\end{itemize}

\subsubsection{Practical Implications for Audio Circuits}
In audio preamplifier design:
\begin{itemize}
  \item \textbf{Low-impedance microphones (150–600 $\Omega$):} favor low-voltage-noise bipolar op-amps (e.g., NE5532, LM4562).
  \item \textbf{High-impedance transducers (piezo, guitar pickups):} favor low-current-noise FET op-amps (e.g., OPA2134, TL072).
  \item \textbf{Tone control and line-level circuits:} both types perform adequately if the impedance network is properly matched.
\end{itemize}

The subjective impact of electronic noise depends on gain and spectral content.  
For instance, $1\,\mu\text{V}_{\mathrm{RMS}}$ of input noise in a 40\,dB preamp becomes $100\,\mu\text{V}$ at the output—corresponding to about $-80\,\text{dBV}$, near the threshold of audibility in quiet studio environments.

\subsubsection{Design Optimization Strategies}
To achieve the lowest practical noise floor:
\begin{enumerate}
  \item Select op-amps whose noise characteristics match the expected source impedance.
  \item Keep resistors in the signal path as low as practical to reduce thermal noise.
  \item Minimize input bias currents and leakage paths to avoid DC offsets that can increase noise coupling.
  \item Limit the system bandwidth to what is strictly necessary (e.g., 20\,Hz–20\,kHz) using well-designed filters.
  \item Pay attention to grounding, shielding, and layout symmetry to avoid additional induced noise.
\end{enumerate}


